\documentclass[UTF8]{ctexart}

\usepackage[a4paper,margin=2.2cm]{geometry}
\usepackage{amsmath}
\usepackage{booktabs}
\usepackage{hyperref}
\usepackage{enumitem}

\title{CS336 Assignment 5 Alignment 总结报告}
\author{}
\date{2026 年 6 月 5 日}

\begin{document}
\maketitle

\section{总体成果}

本次完成的是 Assignment 5 的 Alignment / Reasoning RL 本地实现部分。核心工作是把
\texttt{tests/adapters.py} 中的接口接到真实实现，新增
\texttt{cs336\_alignment/alignment.py}，覆盖 GRPO 训练所需的 tokenization、response
log-prob、reward、advantage normalization、policy-gradient loss、microbatch
aggregation、single train step，以及 supplement 中的 packed SFT dataset、MMLU/GSM8K
解析和 DPO loss。

最终本地验证结果如下：

\begin{center}
\begin{tabular}{lll}
\toprule
验证命令 & 结果 & 说明 \\
\midrule
\texttt{python -m pytest -q} & 26 passed & 本地 Python 环境全量测试 \\
\texttt{UV\_CACHE\_DIR=.uv-cache uv run pytest -q} & 26 passed & 课程 uv/Python 3.12 环境全量测试 \\
\texttt{UV\_CACHE\_DIR=.uv-cache bash test\_and\_make\_submission.sh} & 19 passed & 提交脚本要求的 GRPO 测试 \\
\bottomrule
\end{tabular}
\end{center}

同时生成了 \texttt{test\_results.xml} 和 \texttt{code.zip}。打包脚本已修正为不把
\texttt{.uv-cache} 和 \texttt{.venv} 放入提交包。由于当前机器没有可用的大模型训练资源，
也没有课程共享 Modal/SUNET 环境，本次没有实际运行 OLMo/Llama 的 B200 级训练实验；报告中的
``实验结果''主要来自课程单元测试、snapshot 数值校验和打包脚本验证。

\section{阶段一：阅读与接口定位}

\textbf{任务是什么：} 阅读 Assignment 5 讲义、README、tests 和 adapter，确定需要实现的接口。
README 指明需要完成 \texttt{tests/adapters.py}，必需测试为 GRPO；目录中还包含 supplement
的 data、metrics 和 DPO 测试。

\textbf{用到的大模型知识：} 本阶段主要梳理 alignment pipeline：prompting baseline、GRPO
作为 reasoning RL 算法、SFT instruction tuning、DPO preference optimization，以及评测解析。

\textbf{实验结果：} 初始运行 \texttt{python -m pytest -q} 得到 26 个失败，全部来自
\texttt{NotImplementedError}，说明环境和 fixture 可用，失败原因是接口未实现。

\section{阶段二：Prompt/Output Tokenization 与 Log Prob}

\textbf{任务是什么：} 实现 prompt 和 output 分开 tokenize 后拼接，并构造 causal LM 训练所需的
\texttt{input\_ids}、右移一位的 \texttt{labels} 和只覆盖 response label 的
\texttt{response\_mask}。随后根据模型 logits 计算每个 label token 的条件 log probability 和可选
entropy。

\textbf{用到的大模型知识：} causal LM 的训练目标是
\(\log p(x_t \mid x_{<t})\)。因此模型输入是序列前缀，label 是右移 token；RL 只优化模型生成的
response，不应把 prompt token 纳入 policy-gradient loss。

\textbf{关键代码：}
\begin{verbatim}
input_ids[row, :length] = torch.tensor(full_ids[:max_len])
labels[row, :len(row_labels)] = torch.tensor(full_ids[1:max_len + 1])
response_mask[row, :len(row_labels)] = torch.tensor(
    [(label_position + 1) >= prompt_len
     for label_position in range(len(row_labels))]
)

log_probs_all = F.log_softmax(logits, dim=-1)
log_probs = torch.gather(log_probs_all, dim=-1,
                         index=labels.unsqueeze(-1)).squeeze(-1)
\end{verbatim}

\textbf{实验结果：} \texttt{test\_tokenize\_prompt\_and\_output} 和
\texttt{test\_get\_response\_log\_probs} 通过。第一次实现时短序列 padding 对齐错误，修正为
\texttt{full\_ids[:max\_len]} 与 \texttt{full\_ids[1:max\_len+1]} 后通过 snapshot。

\section{阶段三：Rollout Reward 与 Group Advantage}

\textbf{任务是什么：} 对 rollout response 调用 reward function，返回 raw reward 和
metadata；再按 group size 做 GRPO/Dr.GRPO/MaxRL 风格的 advantage normalization。

\textbf{用到的大模型知识：} GRPO 使用同一个 prompt 下多个 rollout 的组内相对奖励作为 advantage，
通过组均值 baseline 降低方差；不同算法变体改变 baseline、normalizer 或 loss 聚合方式。

\textbf{关键代码：}
\begin{verbatim}
grouped = flat_rewards.reshape(-1, group_size)
centered = grouped - grouped.mean(dim=1, keepdim=True)
advantages = centered / (grouped.std(dim=1, keepdim=True) + advantage_eps)
\end{verbatim}

\textbf{实验结果：} \texttt{compute\_rollout\_rewards}、
\texttt{compute\_group\_normalized\_rewards\_grpo}、
\texttt{compute\_group\_normalized\_rewards\_drgrpo} 和
\texttt{compute\_group\_normalized\_rewards\_maxrl} 全部通过。

\section{阶段四：Policy Gradient Loss 与 GRPO Train Step}

\textbf{任务是什么：} 实现 on-policy loss、off-policy importance reweighting、GRPO/PPO token-level
clipping、GSPO sequence-level clipping，并实现一个完整 single-batch GRPO 训练步。

\textbf{用到的大模型知识：} policy gradient 最大化
\(A_t \log \pi_\theta(a_t|s_t)\)，实现中以 negative objective 作为 loss。off-policy 场景需要
importance ratio \(\exp(\log\pi_\theta-\log\pi_{\text{old}})\)；clip objective 用于控制策略更新幅度。
GSPO 把 token ratio 聚合为 sequence-level geometric ratio。

\textbf{关键代码：}
\begin{verbatim}
ratio = torch.exp(policy_log_probs - old_log_probs)
clipped_ratio = torch.clamp(ratio, 1.0 - cliprange, 1.0 + cliprange)
objective = torch.minimum(advantages * ratio, advantages * clipped_ratio)
loss = -objective

sequence_log_ratio = (log_ratio * response_mask).sum(dim=1, keepdim=True) / denom
sequence_ratio = torch.exp(sequence_log_ratio)
\end{verbatim}

\textbf{实验结果：} on-policy、noclip off-policy、GRPO clip、GSPO clip、sequence normalization、
constant normalization，以及标准/变体 train step 测试均通过。全量 GRPO 测试在提交脚本中为
19 passed。

\section{阶段五：SFT Packed Dataset 与 Batching}

\textbf{任务是什么：} 实现 instruction tuning 的 packed dataset，把 Alpaca 格式样本 tokenize 后
连续拼接，并切成固定长度 chunk；同时实现 dataloader batching。

\textbf{用到的大模型知识：} SFT 仍是 causal LM next-token prediction。packing 可以减少 padding
浪费，提高 token 利用率；样本之间需要 EOS 边界，避免模型把两个独立 instruction-response 样本误看作
连续语义。

\textbf{关键代码：}
\begin{verbatim}
text = prompt_template.format(
    instruction=row["prompt"], response=row["response"]
).rstrip()
token_ids.extend(tokenizer.encode(text, add_special_tokens=True))
token_ids.append(tokenizer.eos_token_id)
\end{verbatim}

\textbf{实验结果：} \texttt{test\_packed\_sft\_dataset} 与
\texttt{test\_iterate\_batches} 通过。第一次实现漏掉 EOS，snapshot 中样本边界不匹配；追加 EOS 后与
预期 tokenized sample 对齐。

\section{阶段六：MMLU/GSM8K 解析}

\textbf{任务是什么：} 从模型输出中解析 MMLU 的 A/B/C/D 选项，以及 GSM8K 的最终数字答案。

\textbf{用到的大模型知识：} 自动评测需要把自由文本 generation 映射为结构化答案。解析函数应尽量稳健：
MMLU 优先匹配 ``answer is''、``answer:''、``option'' 等模式；GSM8K 取输出中最后一个数字作为答案。

\textbf{关键代码：}
\begin{verbatim}
match = re.search(r"\b(?:answer\s+is|answer:|option)\s*([ABCD])\b",
                  model_output, flags=re.IGNORECASE)
matches = re.findall(r"[-+]?(?:\d[\d,]*)(?:\.\d+)?", model_output)
\end{verbatim}

\textbf{实验结果：} \texttt{test\_parse\_mmlu\_response}、
\texttt{test\_parse\_mmlu\_response\_unknown}、
\texttt{test\_parse\_gsm8k\_response} 和
\texttt{test\_parse\_gsm8k\_response\_unknown} 全部通过。

\section{阶段七：DPO Loss}

\textbf{任务是什么：} 给定 policy model、reference model、prompt、chosen response 和 rejected
response，计算单样本 DPO loss。

\textbf{用到的大模型知识：} DPO 直接优化 preference pair，不显式训练 reward model。其核心是比较
policy 相对 reference 的 chosen/rejected log-likelihood gap：
\[
\mathcal{L}_{DPO}
= -\log\sigma\left(\beta\left[
(\log\pi_\theta(y_w|x)-\log\pi_\theta(y_l|x))
-(\log\pi_{ref}(y_w|x)-\log\pi_{ref}(y_l|x))
\right]\right).
\]

\textbf{关键代码：}
\begin{verbatim}
dpo_prompt = prompt_template.format(instruction=prompt, response="")
chosen_response = response_chosen.rstrip() + tokenizer.eos_token
rejected_response = response_rejected.rstrip() + tokenizer.eos_token
preference_logit = beta * (
    (chosen_logp - rejected_logp) - (chosen_ref_logp - rejected_ref_logp)
)
return -F.logsigmoid(preference_logit)
\end{verbatim}

\textbf{实验结果：} \texttt{test\_per\_instance\_dpo\_loss} 通过。调试中发现裸 prompt 得到
loss \(1.1436\)，不符合期望；使用完整 Alpaca 模板并把 EOS 计入 response 后得到
\(0.91038\)，与单测期望 \(0.9104\) 对齐。

\section{失败尝试记录}

\begin{enumerate}[label=\arabic*.]
\item 基线运行 \texttt{python -m pytest -q}：26 failed，全部是 adapter 未实现。
\item 首次实现后运行全量测试：23 passed, 3 failed。失败点分别是 tokenization padding/shift、
packed SFT EOS 边界、DPO prompt 模板。
\item 第一次运行 \texttt{bash test\_and\_make\_submission.sh}：脚本继续打包，但 \texttt{uv run}
因默认缓存目录 \texttt{/home/ldz/.cache/uv} 只读而失败。
\item 使用本地 \texttt{UV\_CACHE\_DIR=.uv-cache} 后，\texttt{uv} 又因网络沙箱无法下载
\texttt{flash-attn} wheel 元数据而失败。联网授权后通过。
\item 为运行 \texttt{uv} 临时生成的 \texttt{.venv} 和 \texttt{.uv-cache} 占用大量空间；验证和打包完成后已删除。
\end{enumerate}

\section{最终状态}

最终代码层面完成了 assignment5 本地可测的全部接口，包含 required GRPO 和 supplement
中的 data、metrics、DPO。最终产物包括：

\begin{itemize}
\item \texttt{cs336\_alignment/alignment.py}: 核心实现。
\item \texttt{tests/adapters.py}: adapter 接口接入实现。
\item \texttt{assignment5\_work\_log.md}: 失败尝试与验证记录。
\item \texttt{assignment5\_report.tex}: 本总结报告。
\item \texttt{test\_results.xml}: 提交脚本生成的测试报告。
\item \texttt{code.zip}: 最终打包文件。
\end{itemize}

受限于本机没有可用 B200/Modal/SUNET 课程资源，本次没有执行讲义中需要大模型 GPU 训练的
OLMo-2/Llama-3.1 实验、AlpacaEval 评测和安全红队实验；这些属于环境资源缺口，而不是本地接口实现缺口。

\end{document}
